\documentclass[12pt,a4paper]{article}
\usepackage[margin=2.3cm]{geometry}
\usepackage{fontspec}
\setmainfont{Times New Roman}
\usepackage{setspace}
\singlespacing
\usepackage{titlesec}
\titleformat{\section}{\normalfont\large\bfseries}{}{0em}{}
\titleformat{\subsection}{\normalfont\normalsize\bfseries}{}{0em}{}
\titlespacing{\section}{0pt}{10pt}{4pt}
\titlespacing{\subsection}{0pt}{8pt}{2pt}
\usepackage{hyperref}
\hypersetup{hidelinks}
\setlength{\parindent}{0pt}
\setlength{\parskip}{4pt}
\newenvironment{refs}{\begin{list}{}{\setlength{\leftmargin}{1.5em}\setlength{\itemindent}{-1.5em}\setlength{\itemsep}{1pt}}}{\end{list}}

\begin{document}

{\large\bfseries BE451 Degree Project in Finance --- Thesis Synopsis, Fall 2026}

\textbf{Students:} [Full name (student ID)], [Full name (student ID)] \\
\textbf{Preferred tutor:} [Alvin Chen / no preference]

\vspace{4pt}\hrule\vspace{4pt}

\section{Topic 1: Do Taxes Depress or Stimulate Corporate Borrowing? Evidence from Sweden's 2019 Tax Cut and Interest Deduction Limit}

\subsection{(i) Tentative research question}

\emph{When Sweden capped the deductibility of interest at 30 percent of EBITDA in 2019 and simultaneously cut the corporate tax rate, did the firms whose tax shield was curtailed reduce leverage and investment, as the trade-off theory predicts, or did the lower rate raise their borrowing capacity, as the productivity channel predicts---and where in the population does the cap actually bind?}

The reform (prop.\ 2017/18:245, in force 1 January 2019) contains two treatments. The statutory rate fell from 22 to 21.4 percent in 2019 and to 20.6 percent in 2021, lowering the interest tax shield for every firm; simultaneously, net interest became deductible only up to 30 percent of taxable EBITDA, with a simplification rule allowing SEK 5 million per group of associated companies regardless. Trade-off theory (Heider and Ljungqvist, 2015) predicts lower leverage from both changes. Sanati and Beyhaghi (2025) show that losing the deduction lowers debt and investment, and Lian and Ma (2021) show that most corporate borrowing is constrained by EBITDA-based limits of exactly the type the cap introduces. Ivanov, Pettit, and Whited (2026) find the opposite for small US private firms: after tax cuts, default moves further away, credit spreads fall, and leverage rises. The cap is our primary test (a sharp bound-versus-unbound split and a kink for bunching); the rate cut is secondary, since 1.4 points implies a small effect. Sweden lets us run both tests in one population of all limited companies: the cap binds only where net interest exceeds both SEK 5 million and 30 percent of EBITDA, giving a clear split between bound and unbound firms and a kink at SEK 5 million. A null on both margins is itself informative: Swedish private-firm capital structure would then not be tax-driven at US elasticities.

\subsection{(ii) Potential data sources and empirical methods}

\textbf{Data.} Serrano (Swedish House of Finance), annual accounts for all Swedish limited companies 2013--2024: total and interest-bearing debt, interest expense, tax expense, EBITDA, investment, dividends, group parent and intra-group balances. Access through SWAMID [verified on (date); N $\approx$ (number) firm-years]. Treatment variables: (a) cap: net interest in 2017--2018 above versus below SEK 5 million and 30 percent of EBITDA, aggregated within groups from Serrano's group identifiers; (b) rate cut: pre-reform tax payments over assets, loss-making firms untreated; (c) rental-housing firms, which received a 12 percent primary deduction, as a third treatment. The cap binds for a few thousand firms, mainly real estate and PE-owned groups.

\textbf{Methods.} Primary test (cap): difference-in-differences for firms bound by the cap in 2017--2018 (net interest above both SEK 5 million and 30 percent of EBITDA) versus unbound firms of similar size, leverage and industry, on leverage, new borrowing, interest expense over debt and investment, with firm and industry-by-year fixed effects; a bunching analysis at the SEK 5 million kink in group net interest (the marginal deduction falls from full to zero at the kink for firms whose EBITDA-based allowance is below SEK 5 million); and a separate treatment of real-estate firms, which received the primary deduction in the same law. Secondary test (rate cut): difference-in-differences on pre-reform tax exposure interacted with post-2019 and post-2021 indicators. Identification rests on the law fixing the kink and the rate path: within industry and year, firms with different pre-reform exposure faced the same macro conditions. Threats and mitigations: manipulation at the kink (bunching is a finding; treatment assigned on pre-reform net interest); intra-group reallocation (group aggregation, following Brok); the 2022 rate cycle (industry-by-year effects, 2016--2018 placebo); 2019 hybrid and CFC rules (exclude large multinationals). Benchmark: the Heider--Ljungqvist elasticity implies about 0.6 percentage points of leverage for the 1.4 point cut; with population data the minimum detectable effect is well below that. Three-month plan: bound-versus-unbound and bunching first, then the group channel and the rate cut. Fallback: cap analysis only.

\subsection{(iii) Short literature review}

\textbf{Ivanov, Pettit, and Whited (2026, Review of Economic Studies)} use US state tax changes and private-firm data and find that leverage rises after tax cuts because lower taxes raise the productivity of capital and lower credit spreads. We replicate their tax-cut test with the 2019 and 2021 rate steps and extend it by adding the simultaneous deduction cap, which their setting lacks. \textbf{Sanati and Beyhaghi (2025, Review of Financial Studies)} exploit the USD 25 million exemption of the 2017 US reform in a regression discontinuity and find that losing the tax benefit of debt lowers debt and investment. We adapt their exemption-threshold design to Sweden's simplification rule, which is a kink rather than a discontinuity, using bunching at the kink and a bound-versus-unbound comparison at the group level. \textbf{Heider and Ljungqvist (2015, Journal of Financial Economics)} estimate that a one percentage point tax increase raises leverage by about 0.4 percentage points; we use this as the benchmark. \textbf{Lian and Ma (2021, Quarterly Journal of Economics)} show that 80 percent of US corporate debt is constrained by cash-flow-based covenants tied to EBITDA; the Swedish cap converts a tax rule into exactly such a constraint for bound firms, which lets us test their prediction that investment responds to EBITDA-linked limits. \textbf{Brok (2026, Journal of Financial and Quantitative Analysis)} shows that tax avoidance acts as a non-debt tax shield that crowds out the debt tax shield, so the leverage response to taxes is weaker for high-avoidance firms; \textbf{Brok (2024, Review of Finance)} shows that corporate group structure creates additional tax incentives for leverage because parents can deduct subsidiary losses. Both imply that tax effects on debt must be measured at the group level, which Serrano's group identifiers allow. Relative to these papers, we replicate both the cap and the tax-cut designs in one reform and one population, testing the trade-off and productivity channels against each other with a cap that binds at the group level. Our contribution is threefold: the first evidence on Sweden's 2019 reform for the full private-firm population; a horse race between two channels with opposite signs; and a hand-built group-level net-interest variable that assigns the kink correctly.

\subsection*{References}
\begin{refs}
\item Brok, Peter, 2024, Debt and taxes: The role of corporate group structure, \emph{Review of Finance} 28, 2051--2082.
\item Brok, Peter, 2026, Debt and taxes: The role of tax avoidance, \emph{Journal of Financial and Quantitative Analysis} 61, 2007--2032.
\item Heider, Florian, and Alexander Ljungqvist, 2015, As certain as debt and taxes: Estimating the tax sensitivity of leverage from state tax changes, \emph{Journal of Financial Economics} 118, 684--712.
\item Lian, Chen, and Yueran Ma, 2021, Anatomy of corporate borrowing constraints, \emph{Quarterly Journal of Economics} 136, 229--291.
\item Ivanov, Ivan T., Luke Pettit, and Toni M. Whited, 2026, Taxes depress corporate borrowing: Evidence from private firms, \emph{Review of Economic Studies} 93, 2715--2749.
\item Sanati, Ali, and Mehdi Beyhaghi, 2025, How does removing the tax benefits of debt affect firms? Evidence from the 2017 U.S. tax reform, \emph{Review of Financial Studies}, forthcoming, doi:10.1093/rfs/hhaf041.
\end{refs}

\newpage

\section{Topic 2: Investment and Payout Policy under a Threatened Profit Cap: Evidence from Sweden's Private Welfare Firms}

\subsection{(i) Tentative research question}

\emph{How do firms adjust investment, payout policy and ownership when the return on capital retained in the firm may soon be capped by law, and does the adjustment reverse when the threat is withdrawn?}

Political uncertainty is known to depress corporate investment: firms delay irreversible projects until policy is resolved (Julio and Yook, 2012; Jens, 2017). A threatened profit cap adds a second, sharper channel. If retained capital will earn at most a regulated return, its value inside the firm falls, so owners should accelerate dividends and group contributions while the cap is pending, and reverse afterwards. The competing hypothesis is that owners assigned the proposal a low probability, so nothing happened. Sweden's tax-financed welfare sector provides the setting: every private provider of schooling, social care and personal assistance faced a credible cap proposal between the Reepalu inquiry (8 November 2016), the government bill (19 March 2018) and the Riksdag's rejection (7 June 2018), and the threat returned with a directive in June 2022, a bill in June 2026 and the election of 13 September 2026. Health-care providers, which the bill excluded, form a natural control group. To our knowledge no study examines the firm-level financial response. A null is informative: it measures how much a politically salient but legislatively weak threat costs in investment and payouts.

\subsection{(ii) Potential data sources and empirical methods}

\textbf{Data.} Serrano, annual accounts for all Swedish limited companies 2012--2024: investment, dividends, group contributions, cash, debt, employees, entry, exit and group parent. Treated firms: private firms in SNI 85 (education) and 87--88 (residential care, social work, personal assistance). Placebo: SNI 86 (health care), which shares public financing and demand structure but was excluded from the bill. Controls: matched private service firms. Ownership type (private equity or not) hand-coded from group parents and annual reports for the largest welfare groups. Listed welfare firms (Attendo, Humana, AcadeMedia, Ambea) from LSEG or FinBas for an illustrative event study. Access: all course resources [Serrano verified on (date); N $\approx$ (number) treated firms].

\textbf{Methods.} Event-study difference-in-differences with firm and industry-by-year fixed effects around 8 November 2016, 19 March 2018 and 7 June 2018; outcomes investment over assets, dividends over equity, group contributions, cash, employment growth, exit and ownership change. Identification: the treated set is defined by the bill's own scope, and SNI 86 is the within-sector placebo; matched service firms absorb aggregate shocks. Mechanism tests: private-equity-owned versus other treated firms (shorter horizon, higher leverage), and firms with pre-threat profits above versus below the proposed cap (where the cap would bind). Robustness: pre-2016 pseudo-events, exclusion of asylum accommodation (SNI 87.9), alternative matching, and the 2022--2026 school episode as a second treatment. Threats and mitigations: owner and management screening rules from 1 January 2019 (main window 2016--2018; SNI 85 separated from 87--88); the rate shock in the second episode (matching and placebo); the small listed sample (event study illustrative only). Benchmark: Julio and Yook report a 4.8 percent investment reduction in election years; with several thousand treated firms our minimum detectable effect is well below that. Three-month plan: the 2014--2020 difference-in-differences on investment and payouts with the SNI 86 placebo first; then the mechanism splits, the second episode and the event study.

\subsection{(iii) Short literature review}

\textbf{Julio and Yook (2012, Journal of Finance)} document across 48 countries that firms reduce investment by 4.8 percent on average in national election years and recover afterwards, consistent with delaying irreversible investment until uncertainty is resolved. We replicate their event-window design on a single, precisely dated policy threat within one country and extend it by adding payout and ownership outcomes. \textbf{Jens (2017, Journal of Financial Economics)} shows the effect is causal using staggered US gubernatorial elections and placebo tests; we replicate the placebo logic with health-care providers excluded from the bill as a within-sector control. \textbf{Kelly, P\'astor, and Veronesi (2016, Journal of Finance)} show that political uncertainty is priced in option markets and that its price depends on the probability of a policy change; we use their framework for the payout prediction and for the illustrative event study on the four listed welfare firms. \textbf{Bonaime, Gulen, and Ion (2018, Journal of Financial Economics)} show that policy uncertainty reduces merger activity; we use their result for the ownership-change outcome, including sales to private equity. Relative to Julio and Yook (2012) and Jens (2017), we replicate the political-uncertainty investment design on a legislative threat that hit an identifiable set of private firms, allowing us to test the investment channel and the payout channel, which make opposite predictions for dividends, with a placebo built into the bill itself. Our contribution is threefold: the first firm-level evidence on the financial response to a sector-specific profit-cap threat, using the full population of affected firms; a test of the real-option channel against the payout channel with private-equity ownership as a mechanism split; and hand-coded legislative dates, sector scope and ownership of the largest welfare groups.

\subsection*{References}
\begin{refs}
\item Bonaime, Alice, Huseyin Gulen, and Mihai Ion, 2018, Does policy uncertainty affect mergers and acquisitions?, \emph{Journal of Financial Economics} 129, 531--558.
\item Jens, Candace E., 2017, Political uncertainty and investment: Causal evidence from U.S. gubernatorial elections, \emph{Journal of Financial Economics} 124, 563--579.
\item Julio, Brandon, and Youngsuk Yook, 2012, Political uncertainty and corporate investment cycles, \emph{Journal of Finance} 67, 45--83.
\item Kelly, Bryan, \v{L}ubo\v{s} P\'astor, and Pietro Veronesi, 2016, The price of political uncertainty: Theory and evidence from the option market, \emph{Journal of Finance} 71, 2417--2480.
\end{refs}

\end{document}
